\scanpage{4}\mspage{\handwritten{3}}

premier. Comme application, Quillen obtient une interprétation \unsure{faisceautique} de la cohomologie d'un ensemble semi-simplicial à coéfficients dans \unsure{un} système \inserted{\uline{local}} contravariant \deleted{\unsure{de}} \inserted{$F$} (défini classiquement par le complexe \unsure{cosimplicial} des $C^n(F)=\handwritten{\coprod}_{x\in X_n}F(x)$)\margin{\handwritten{$\prod$ ?}} : on considère le système local contravariant défini par $F$, on l'interprète comme un faisceau sur \inserted{$T(X)$ i.e. objet de} $(\mathrm{Ssimpl})_{/X}$, et on prend sa cohomologie. - Cependant, quand $F$ est un système de coéfficien\unsure{bs|ts|s} covariant pas nécessairement local, on n'a tjrs pas d'interp\unsure{éé|ré}tation de \unsure{ses} groupes de cohomologie classiques en termes faisceautiques; ni, lorsque \unsure{$F$} est contravariant, de son homologie, ou inversement de sa cohomologie faisceautique en termes classiques.\margin{\handwritten{Si ! \tikz[scale=.18,baseline]{\draw (.5,1)--(0,1)--(0,-1)--(.5,-1);}}}

\margin{\handwritten{\(\Vert\)}}% parallel pencil rules beside the Quillen assertion
A propos de la notion de foncteur qui est un homotopisme. Quillen montre qu\deleted{e}'un tel foncteur $F:C\handwritten{\to}C'$ induit une équivalence entre la sous catégorie triangulée \inserted{$D^b_{lc}(C')$} de la catégorie \inserted{dérivée \inserted{bornée} de celle \deleted{des}} des faisceaux abéliens sur $C'$, dont les faisceaux de cohomologie sont des systèmes locaux, et la catégorie \unsure{analogue} pour $C$; \deleted{\unsure{il s'en}} et réciproquement. \deleted{Pour la partie directe} On peut dans cet énoncé introduire aussi n'importe quel anneau de base (à \unsure{condition} de le supposer $\handwritten{\ne}0$ dans le cas de la réciproque); \deleted{constant} la partie \unsure{directe} vaut aussi avec un anneau de boefficients pas néc. constant, mais \unsure{constant} tordu. Je pense que ce résultat (facile) doit pouvoir se généraliser \unsure{ainsi} Soit $f:X\quad X'$ un morphisme de topos qui soit tel que pour tout \deleted{\unsure{système}} faisceau loc. const. sur $X'$, $f$ induise un isomorphisme sur les \unsure{cohomologies} (avec cas non commutatif inclus). Supposons que \deleted{\unsure{le topos $X$|$X$ (resp. $X'$)}} \inserted{$X$\deleted{\unsure{$'$}} et $\handwritten{X'}$} soit \uline{\unsure{localement} homotopiquement trivial}, i.e. que \deleted{tout} pour tout entier $n\inserted{\handwritten{\geq}1}$, tout objet $U$ \deleted{\unsure{de $X'$}} ait un \deleted{revêtement}\inserted{recouvrement} par des $U_i\handwritten{\to}U$, tels que a) tout système local sur \unsure{$U$} devient constant sur $U_i$, et toute section sur $U$ devient constante sur $U_i$ et b) pour tout \deleted{système} groupe abélien $G$, les $H^{\deleted{\unsure{i}}j}(U,G)\handwritten{\to}H^j(U_i,G)$ sont nuls pour $1\handwritten{\leq}j\handwritten{\unsure{<|\leq}}n$.\margin{\handwritten{\tikz[scale=.16,baseline]{\draw (0,0) circle (.28) (.15,.2)--(1.25,1.75);}\ \(\vert\!\vert\!\vert\)}}% callout and three vertical strokes beside the condition
\margin{\handwritten{\tikz[scale=.16,baseline]{\draw (0,0) circle (.28) (.15,.2)--(1.25,1.75);}\ Attention, cette condition n'est typiquement pas satisfaite par les schémas en leur \unsure{topologie} \unsure{étale} (mais bien par \unsure{schémas} \unsure{noeth.|loc. noeth.} avec top. Zariski).}} Alors le foncteur $D^b_{lc}(X')\handwritten{\to}D^b_{lc}(X)$ induit par $f$ est

% archive sheet number: 3
